\begin{exsol}{1}{\MR{} Food for Thought 1.1}{18}{Foundations}{A1}
\asked Say which decisions in a fast-food chain are taken by headquarters and which by the outlet
managers, explain why, and judge whether splitting the problems into coordination and motivation
is a useful way to analyse the business.

\method The organizing pair of the course: coordination (making plans fit, an information problem)
versus motivation (getting people to carry them out, an interest problem) (\MR{} Ch.~1, p.~16;
Ch.~2, ``The Tasks of Coordination and Motivation'', pp.~25--27). Authority goes where the
information is, and incentives and delegated authority are complements (p.~17). Design variables
are set centrally and details are left to local managers (Ch.~4, pp.~91, 114--115). Routines handle
recurrent design problems (p.~92). Scale and scope economies (pp.~106--107).

\solution
\subparagraph{What the centre decides, and why.} The menu and recipes, product specifications and
portion sizes, food-safety and cleanliness standards, the operating manual, store layout and
equipment, brand and national advertising, approved suppliers and central purchasing, and the
launch of new products. These are \emph{design variables} in the sense of p.~91. First, there is
strong a priori information about the optimum: the chain's value is the promise that the product
is the same in every outlet, so every outlet must match the standard. Second, a misfit is costly:
one outlet that serves a different or unsafe product damages the brand everywhere. Many of these
choices also carry scale economies (purchasing, advertising, product development; pp.~106--107)
that no single outlet can capture. New products have \emph{innovation attributes} (p.~92): testing a
new item needs information and resources beyond any one outlet's experience.

\subparagraph{What the outlet decides, and why.} Crew hiring, training on the floor and
supervision, shift schedules, the quantities ordered from approved suppliers, day-to-day
maintenance, handling complaints and some local promotions. The information these need is local
and changes fast: Friday's lunch rush, the weather, a local event, which worker is reliable. Sending
it to headquarters would be slow and costly, and the information would be stale on arrival
(Ch.~2, ``The Need for Information'', p.~26). Following p.~17, authority goes where the
information is. The operating manual turns many recurring choices into \emph{routines} (p.~92),
so the manager solves them locally without asking the centre. In the language of p.~114, the
central decisions set the parameters within which the local ones are made.

\subparagraph{Coordination and motivation.} The \emph{coordination} problem is making thousands of
kitchens produce an identical product at a predictable quality. It is solved mostly by central
standards and routines. The \emph{motivation} problem is that the manager's effort on cleanliness,
speed and courtesy is hard to observe, and most of the gain from that effort goes to the brand, not
to the outlet. An outlet serving travellers who rarely return bears little of the cost of shirking:
the damage falls on all other outlets. The tools are an ownership stake in the outlet's profit
(the franchise form makes the manager a residual claimant), plus inspections, standards backed by
termination and required input sources. Delegating local decisions pays only because the decision
maker shares in the result (p.~17: incentives and delegation are complements).

\subparagraph{Is the split useful?} Yes, as a lens, but not as a way to sort decisions into two
piles. Most decisions have both sides. Hygiene standards are centralized partly for coordination
(uniformity) and partly for motivation: left alone, an outlet would under-invest in a brand it
shares with others. Profit-based pay for a franchisee also creates a new motivation problem, since
cutting costs below the standard raises its own profit at the brand's expense; the centre then
monitors to protect the coordinated design. As \MR{} put it, motivation problems ``can influence
the effectiveness of coordination mechanisms'' (p.~114). The good answer uses the split to ask two
questions of each decision: who has the information, and whose interests does it touch?

\interp This is the Aula~1 lesson on a small scale: the chain is a coherent system in which central
standards, routines, local discretion and residual-claimant incentives support one another, the
way GM's divisions, accounting and central office did (p.~17).

\verify Limit cases. With no brand externality (independent restaurants), nearly everything would
be decided locally and coordination through a central standard would lose its point. With
perfectly observable effort, outlet managers could be salaried and closely directed, and ownership
incentives would matter less. The observed split sits between the two, as the theory predicts.

\answer{The centre sets the design variables --- menu, recipes, standards, layout, suppliers,
branding, new products --- because the optimum is known to require uniformity, misfits damage the
whole brand, and they carry scale economies. Outlets decide staffing, scheduling, ordering and
service, where information is local and fast-changing. The coordination/motivation split is a
useful lens, but each decision has both aspects: standards are also a defence against
free-riding on the brand, and outlet incentives must be paired with monitoring.}
\end{exsol}
